\section{Three Empirical Anchors}

The empirical role of the indexed sources is narrow and disciplined:
\begin{enumerate}[nosep]
  \item \textbf{Signatures of conscious state.} The 2024 Communications
        Biology source is used for one strong claim only: conscious conditions
        are associated with robust and generalizable dynamical
        structure-function signatures rather than being describable only in
        loose psychological language.
  \item \textbf{Connectivity and stability.} The 2025 Communications Biology
        source is used for a second strong claim: global conscious states are
        related to structured connectivity, local dynamics, and stability, so
        persistence is not mere noise but an architectural problem.
  \item \textbf{Transition dynamics.} The 2025 Nature Neuroscience sleep
        source is used for a third strong claim: entry into a new regime can
        be abrupt, patterned, and tractable in dynamical language rather than
        remaining purely descriptive.
\end{enumerate}

\begin{discussion}
These sources do not prove the full framework-inertia model, and the chapter
should not pretend otherwise. Their value is more disciplined. They justify
treating state, stability, and transition as serious biological topics. Once
that floor is established, the rest of the manuscript can talk about entry
cost, regime persistence, and redirection without sounding like a private
productivity metaphor detached from the science of organized state changes.
\end{discussion}
